\subsection{Modules of fractions}

The canonical homomorphism \(i_{\mathbf{A}}^{\mathbf{S}} \colon \mathbf{A} \to \mathbf{A}[\mathbf{S}^{-1}]\) defined in no. 1 allows us to consider every \(\mathbf{A}[\mathbf{S}^{-1}]\) -module as an A-module.

PROPOSITION 3. Let A be a ring, S a subset of A, M an A-module, M' the A-module M ⊗\_A A{[}S\^{}\{-1\}{]} and f the canonical A-homomorphism x ↦ x ⊗ 1 of M to M'. Then:

\begin{enumerate}
\def\labelenumi{(\arabic{enumi})}
\item
  For all \(s \in \mathbf{S}\) , the homothety \(z \mapsto sz\) of \(\mathbf{M}'\) is bijective.
\item
  For every A-module N such that, for all \(s \in S\) , the homothety \(y \mapsto sy\) of N is bijective, and every homomorphism u of M to N, there exists a unique homomorphism \(u'\) of \(M'\) to N such that \(u = u' \circ f\) .
\end{enumerate}

In other words, \((M',f)\) is a solution of the universal mapping problem (Set Theory, Chapter IV, §3, no. 1) with the following conditions: the species of structure \(\Sigma\) is that of an A-module in which the homotheties induced by the elements of S are bijective, the morphisms are A-module homomorphisms and the \(\alpha\) -mappings are also A-module homomorphisms.

For every A-module N and all \(a \in \mathbf{A}\) , denote by \(h_a\) the homothety \(y \mapsto ay\) in \(\mathbf{N}\) ; \(a \mapsto h_a\) is then a ring homomorphism from A to \(\operatorname{End}_{\mathbf{A}}(\mathbf{N})\) . To say that \(h_a\) is bijective means that \(h_a\) is an invertible element of \(\operatorname{End}_{\mathbf{A}}(\mathbf{N})\) . Suppose that, for all \(s \in \mathbf{S}\) , \(h_s\) is invertible in \(\operatorname{End}_{\mathbf{A}}(\mathbf{N})\) ; the elements \(h_a\) , where \(a \in \mathbf{A}\) , and the inverses of the elements \(h_s\) , where \(s \in \mathbf{S}\) , then generate in \(\operatorname{End}_{\mathbf{A}}(\mathbf{N})\) a commutative subring B and the homomorphism \(a \mapsto h_a\) from A to B is such that the images of the elements of S are invertible. Then it follows (no. 1, Proposition 1) that there exists a unique homomorphism \(h'\) of \(\mathrm{A}[\mathrm{S}^{-1}]\) to B such that

\[
h ^ {\prime} (a / s) = h _ {a} \left(h _ {s}\right) ^ {- 1};
\]

we know (Algebra, Chapter II, § 1, no. 14) that such a homomorphism defines on N an A{[}S \(^{-1}\) {]}-module structure such that (a/s).y = h \(_{s}^{-1}\) (a.y); the A-module structure derived from this A{[}S \(^{-1}\) {]}-module structure by means of the homomorphism i \(_{A}^{S}\) is precisely the structure given initially.

Conversely, if N is an A{[}S \(^{-1}\) {]}-module and it is considered as an A-module by means of \(i_{\text{A}}^{\text{S}}\) , the homotheties \(y \mapsto sy\) , for \(s \in \text{S}\) , are bijective, for \(y \mapsto (1/s)y\) is the inverse mapping of \(y \mapsto sy\) ; and the A{[}S \(^{-1}\) {]}-module structure on N derived from its A-module structure by the process described above is the A{[}S \(^{-1}\) {]}-module structure given initially. Thus there is a canonical one-to-one correspondence between A{[}S \(^{-1}\) {]}-modules and A-modules in which the homotheties induced by the elements of S are bijective; moreover, if N, N' are two A-modules with this property every A-module homomorphism \(u: \text{N} \to \text{N}'\) is also a homomorphism of the A{[}S \(^{-1}\) {]}-module structures of N and N', as, for all \(y \in \text{N}\) and all \(s \in \text{S}\) , we may write \(u(y) = u(s. ((1/s)y)) = s.u((1/s)y)\) , whence \(u((1/s)y) = (1/s)u(y)\) ; the converse is obvious.

This being so, the statement of Proposition 3 is just the characterization of the module obtained from M by extending the scalars to A{[}S \(^{-1}\) {]}, taking account of the above interpretation (Algebra, Chapter II, § 5, no. 1, Remark 1).

DEFINITION 3. Let A be a ring, S a subset of A, \(\overline{\mathbf{S}}\) the multiplicative subset of A generated by S and M an A-module. Then the module of fractions of M defined by S and denoted by \(\mathbf{M}[\mathbf{S}^{-1}]\) or \(\overline{\mathbf{S}}^{-1}\mathbf{M}\) is the \(\mathbf{A}[\mathbf{S}^{-1}]\) -module \(\mathbf{M} \otimes_{\mathbf{A}} \mathbf{A}[\mathbf{S}^{-1}]\) .

In this chapter we shall usually denote by \(i_{\mathbf{M}}^{\mathbf{S}}\) the canonical mapping \(m \mapsto m \otimes 1\) of \(\mathbf{M}\) to \(\mathbf{M}[\mathbf{S}^{-1}]\) .

Remarks

\begin{enumerate}
\def\labelenumi{(\arabic{enumi})}
\item
  Clearly \(\mathbf{M}[\overline{\mathbf{S}}^{-1}] = \mathbf{M}[\mathbf{S}^{-1}]\) .
\item
  For \(m \in \mathbf{M}\) and \(s \in \overline{\mathbf{S}}\) we also write \(m / s\) for the element \(m \otimes (1 / s)\) of \(\mathbf{M}[\mathbf{S}^{-1}]\) . Every element of \(\mathbf{M}[\mathbf{S}^{-1}]\) is of the form, for such an element is of the form \(\sum_{i} m_i \otimes (a_i / s)\) , where \(m_i \in \mathbf{M}\) , \(a_i \in \mathbf{A}\) , \(s \in \mathbf{S}\) (no. 1, Remark 2), and
\[
m _ {i} \otimes (a _ {i} / s) = (a _ {i} m _ {i}) \otimes (1 / s),
\]

hence \(\sum_{i} m_{i} \otimes (a_{i}/s) = m \otimes (1/s)\) , where \(m = \sum_{i} a_{i} m_{i}\) . Then


\[
(m / s) + (m ^ {\prime} / s ^ {\prime}) = (s ^ {\prime} m + s m ^ {\prime}) / s s ^ {\prime}\tag{4}
\]

\[
(a / s) (m / s ^ {\prime}) = (a m) / (s s ^ {\prime})\tag{5}
\]

where \(m, m' \in \mathbf{M}, a \in \mathbf{A}\) and \(s, s' \in \mathbf{S}\) .

\item
  If S is the complement of a prime ideal p of A, we write \(M_{p}\) instead of \(S^{-1}M\) .
\item
  Let M be an A{[}S \(^{-1}\) {]}-module; if M is considered canonically as an A-module, \(i_{M}^{S}\) is a bijection, for the ordered pair consisting of M and the identity mapping \(l_{M}\) is also trivially a solution of the universal mapping problem solved by M{[}S \(^{-1}\) {]} and \(i_{M}^{S}\) . Then M is identified with M{[}S \(^{-1}\) {]}.
\end{enumerate}

PROPOSITION 4. Let S be a multiplicative subset of A and M an A-module. For m/s = 0 (m ∈ M, s ∈ S), it is necessary and sufficient that there exist s' ∈ S such that s'm = 0.

If \(s' \in S\) is such that \(s'm = 0\) , clearly \(m/s = (s'm)/(ss') = 0\) . Conversely, suppose that \(m/s = 0\) . As \(1/s\) is invertible in \(S^{-1}A\) , \(m/1 = 0\) . For every sub-A-module \(P\) of \(S^{-1}A\) containing 1, we denote by \(\beta(P, m)\) the image of \((m, 1)\) under the canonical mapping of \(M \times P\) to \(M \otimes_A P\) ; then \(\beta(S^{-1}A, m) = 0\) . We know (Algebra, Chapter II, §6, no. 3, Corollary 4 to Proposition 7) that there exists a finitely generated submodule \(P\) of \(S^{-1}A\) containing 1 and such that \(\beta(P, m) = 0\) . For all \(t \in S\) we denote by \(A_t\) the set of \(a/t\) , where \(a \in A\) ; as \(P\) is finitely generated, there exists \(t \in S\) such that \(P \subset A_t\) (no. 1, Remark 2), whence \(\beta(A_t, m) = 0\) . The mapping \(a \mapsto a/t\) from \(A\) to \(A_t\) is surjective; let \(B\) be its kernel. It defines a surjective mapping \(h: M \otimes_A A \to M \otimes_A A_t\) , whose kernel is BM (M being identified with \(M \otimes A\) ); then \(\beta(A_t, m) = h(tm)\) and consequently \(tm\) can be expressed in the form \(\sum_{i=1}^{r} b_i m_i\) , where \(b_i \in B\) , \(m_i \in M (1 \leqslant i \leqslant r)\) . As \(b_i/t = 0\) expressed in the form \(\sum_{i=1}^{n} b_i m_i\) , where \(b_i \in \mathbf{B}\) , \(m_i \in \mathbf{M}\) ( \(1 \leqslant i \leqslant r\) ). As \(b_i / t = 0\) for \(1 \leqslant i \leqslant r\) , there exists \(t' \in \mathbf{S}\) such that \(t'b_i = 0\) for \(1 \leqslant i \leqslant r\) , whence \(t'tm = 0\) , which proves Proposition 4.

COROLLARY 1. For \(m / s = m' / s'\) in \(S^{-1}M\) , it is necessary and sufficient that there exist \(t \in S\) such that \(t(s'm - sm') = 0\) .

\[
(m / s) - (m ^ {\prime} / s ^ {\prime}) = (s ^ {\prime} m - s m ^ {\prime}) / s s ^ {\prime}.
\]

COROLLARY 2. Let M be a finitely generated A-module. For \(S^{-1}M = 0\) , it is necessary and sufficient that there exists \(s \in S\) such that \(sM = 0\) .

Without any conditions on M, clearly the relation \(sM = 0\) for some \(s \in S\) implies \(S^{-1}M = 0\) . Conversely, suppose that \(S^{-1}M = 0\) and let \((m_i)_{1 \leqslant i \leqslant n}\) be a system of generators of M; the \(m_i / 1\) generate the \(S^{-1}A\) -module \(S^{-1}M\) , hence to say that \(S^{-1}M = 0\) amounts to saying that \(m_i / 1 = 0\) for \(1 \leqslant i \leqslant n\) ; by

Proposition 4 there exist \(s_i \in S\) such that \(s_i m_i = 0\) and, taking \(s = s_1 s_2 \ldots s_n\) , \(sm_i = 0\) for all \(i\) and hence \(sM = 0\) .

COROLLARY 3. Let M be a finitely generated A-module. For an ideal a of A to be such that aM = M, it is necessary and sufficient that there exist \(a \in a\) such that \((1 + a)M = 0\) .

Clearly the relation \((1 + a)M = 0\) implies M = aM. To prove the converse we use the following lemma:

LEMMA 1. For every ideal a of A, the set S of elements \(1 + a\) , where \(a \in a\) , is a multiplicative subset of A and the set \(a'\) of elements of \(S^{-1}A\) of the form a/s, where \(a \in a\) and \(s \in S\) , is an ideal contained in the Jacobson radical of \(S^{-1}A\) .

The first assertion is obvious, as well as the fact that \(a'\) is an ideal of \(S^{-1}A\) . On the other hand, \((1/1) + (a/s) = (s + a)/s\) and \(s + a \in S\) for all \(s \in S\) and \(a \in a\) by definition of S; hence \((1/1) + (a/s)\) is invertible in \(S^{-1}A\) for all \(a/s \in a'\) , which completes the proof of the lemma (Algebra, Chapter VIII, §6, no. 3, Theorem 1).

This being so, if we set \(N = S^{-1}M\) , clearly N is a finitely generated \(S^{-1}A\) -module; if aM = M, then \(a'N = N\) and it follows that N = 0 by Nakayama's Lemma (Algebra, Chapter VIII, §6, no. 3, Corollary to Proposition 6); the corollary then follows from Corollary 2.

PROPOSITION 5. Let A, B be two rings, S a multiplicative subset of A, T a multiplicative subset of B and f a homomorphism from A to B such that \(f(S) \subset T\) . Let M be an A-module, N a B-module and u an A-linear mapping from M to N. Then there exists a unique \(S^{-1}A\) -linear mapping \(u'\) from \(S^{-1}M\) to \(T^{-1}N\) such that \(u'(m/1) = u(m)/1\) for all \(m \in M\) .

The mapping \(i_{\mathbf{N}}^{\mathbf{T}} \circ u\) from \(\mathbf{M}\) to \(\mathbf{T}^{-1}\mathbf{N}\) is A-linear. Moreover, if \(s \in \mathbf{S}\) , then \(f(s) \in \mathbf{T}\) , hence the homothety induced by \(s\) on \(\mathbf{T}^{-1}\mathbf{N}\) is bijective. The existence and uniqueness of \(u'\) then follow from Proposition 3. Then, for \(m \in \mathbf{M}\) and \(s \in \mathbf{S}\) ,

\[
u ^ {\prime} (m / s) = u (m) / f (s).\tag{6}
\]

With the same notation, let C be a third ring, U a multiplicative subset of C, g a homomorphism from B to C such that \(g(T) \subset U\) , P a C-module, v a B-linear mapping from N to P and \(v'\) the \(T^{-1}B\) -linear mapping from B to \(U^{-1}P\) associated with v. Then

\[
(v \circ u) ^ {\prime} = v ^ {\prime} \circ u ^ {\prime}\tag{7}
\]

where the left hand side is the A-linear mapping \(S^{-1}M \rightarrow U^{-1}P\) associated with \(v \circ u\) . Similarly, if \(u_{1}\) is a second A-linear mapping from M to N, then

\[
(u + u _ {1}) ^ {\prime} = u ^ {\prime} + u _ {1} ^ {\prime},\tag{8}
\]

the left-hand side being the A-linear mapping \(S^{-1}M \rightarrow T^{-1}N\) associated with \(u + u_{1}\) .

Remark (5). If, in Proposition 5, we take B = A, T = S and \(f = 1_{A}\) , it is easily seen that \(u'\) is just the mapping \(u \otimes 1: M \otimes S^{-1}A \to N \otimes S^{-1}A\) . We shall henceforth denote it by \(S^{-1}u\) ; if S is the complement of a prime ideal p of A, we write \(u_{p}\) instead of \(S^{-1}u\) .

PROPOSITION 6. Let \(f\) be a homomorphism from a ring \(\mathbf{A}\) to a ring \(\mathbf{B}\) and \(\mathbf{S}\) a multiplicative subset of \(\mathbf{A}\) . There exists a unique mapping \(j\) from \((f(\mathbf{S}))^{-1}\mathbf{B}\) to \(\mathbf{S}^{-1}\mathbf{B}\) (where \(\mathbf{B}\) is considered as an A-module by means of \(f\) ) such that \(j(b/f(s)) = b/s\) for all \(b \in \mathbf{B}\) , \(s \in \mathbf{S}\) . If \(f': \mathbf{S}^{-1}\mathbf{A} \to (f(\mathbf{S}))^{-1}\mathbf{B}\) is the ring homomorphism associated with \(f\) (no. 1, Proposition 2), then \(j \circ f' = \mathbf{S}^{-1}f\) . The mapping \(j\) is an isomorphism of the \(\mathbf{S}^{-1}\mathbf{A}\) -module structure on \((f(\mathbf{S}))^{-1}\mathbf{B}\) defined by \(f'\) onto that on \(\mathbf{S}^{-1}\mathbf{B}\) and also of the B-module structure on \((f(\mathbf{S}))^{-1}\mathbf{B}\) onto that on \(\mathbf{S}^{-1}\mathbf{B}\) (resulting from the definition \(\mathbf{S}^{-1}\mathbf{B} = (\mathbf{S}^{-1}\mathbf{A}) \otimes_{\mathbf{A}} \mathbf{B}\) ).

If \(b, b'\) are in B, \(s, s'\) in S, the conditions \(b / s = b' / s'\) and \(b / f(s) = b' / f(s')\) are equivalent, as follows from Corollary 1 to Proposition 4, which establishes the existence of \(j\) and shows that \(j\) is bijective; the uniqueness of \(j\) is obvious. Clearly \(j\) is an additive group isomorphism. If \(a \in A\) , \(b \in B\) , \(s \in S\) , \(t \in S\) , then

\[
(a / s). (b / f (t)) = f ^ {\prime} (a / s) (b / f (t)) = f (a) / f (s) (b / f (t)) = (f (a) b) / f (s t),
\]

from which its follows that \(j\) is \((\mathbf{S}^{-1}\mathbf{A})\) -linear. Clearly \(j \circ f^1 = \mathbf{S}^{-1}f\) . Finally, if \(b \in \mathbf{B}\) , \(b' \in \mathbf{B}\) , \(s \in \mathbf{S}\) , then \(j(b \cdot (b'/f(s)) = j(bb'/f(s)) = bb'/s = b \cdot (b'/s))\) , which proves the last assertion.

The mapping \(j\) of Proposition 6 is called the canonical isomorphism of \((f(\mathbf{S}))^{-1}\mathbf{B}\) onto \(\mathbf{S}^{-1}\mathbf{B}\) . These two sets are in general identified by means of \(f\) ; then \(f' = \mathbf{S}^{-1}f\) , \(i_{\mathbf{B}}^{\mathbf{S}} = i_{\mathbf{B}}^{f(\mathbf{S})}\) .

